\section{[Old appendix by Ulrich] Galois FRI-Binius}

\ulrich{Just notes, not sure if it worth to keep it in final doc} \albert{Nice, thanks!}

Packed commitment scheme for $\mathbb F_{2^m}$ . % in the values-to-uni-coeffs regime.

For simplicity we assume packing across polynomials. 
(The single poly case goes likewise.)
Using an $\mathbb F_2$-basis $\beta_1,\ldots, \beta_{m}$ of $\mathbb F_{2^m}$ we
pack $m$ polys of $2^n$ coefficients %of $2^{m}$ coeffs in $\mathbb F_2$,
\[
b_i(X) = \sum_{j=0}^{2^n - 1} b_{i,j} \cdot X^i \quad \in \mathbb F_2[X]^{< 2^n}, 
\quad i = 1, \ldots, m,
\]
coefficient-wise into a single poly 
\begin{align*}
b(X) &= b_1(X)\cdot \beta_1 + \ldots + b_{m}(X) \cdot\beta_m
\\
&= \sum_{j=0}^{2^n} (b_{1,j}\cdot \beta_1 + \ldots + b_{m,j}\cdot \beta_{m}) \cdot X^j 
\quad \in \mathbb F_{2^m}[X]^{<2^n}.
\end{align*}

The Frobenius automorphism $\phi(z) = z^2$ keeps $\mathbb F_2\subset F$ fixed, and has no other fixed points. 
% (Recall that it is also additive.)
For any point $z\in \mathbb F$, we extract the value of $b_i(z)$ from the values of the packed poly $b(X)$ along the Galois orbit 
\[
z, z^2, z^{2^2}, \ldots, z^{2^{m-1}}
\]
as follows.
%
The Moore matrix
\begin{align*}
M = \begin{pmatrix}
\beta_1 & \beta_2 & \ldots & \beta_m
\\
\beta_1^2 & \beta_2^2 & \ldots & \beta_m^2
\\
\beta_1^{2^2} & \beta_2^{2^2} & \ldots & \beta_m^{2^2}
\\
\vdots
\\
\beta_1^{2^{m-1}} & \beta_2^{2^{m-1}} & \ldots & \beta_m^{2^{m-1}}
\end{pmatrix}
\end{align*}
consists of the Galois orbits of the basis vectors, and is invertible. (See below.)
Using that the component polynomials are over $\mathbb F_2$, we see that
%We write $v_{k} = b(\phi^{-k} z)$, $k=0, \ldots, m-1$, and 
% \[
% v_k = \phi^{-k}\left(b_1(z) \cdot \phi^{k}\beta_1 + \ldots + b_m(z) \cdot \phi^{k}\beta_m \right),
% \]
% since $b_i(\phi^{-k} z) = \phi^{-k} b_i(z)$, in other words
\[
\phi^{k} \left(b(\phi^{-k} z)\right) = b_1(z) \cdot \phi^{k}\beta_1 + \ldots + b_m(z) \cdot \phi^{k}\beta_m,
\]
for any $k$. 
To extract the $b_1(z)$ term, we look for weights $c_1,\ldots, c_m\in F$ which annihilate the orbits of the other $\beta_2,\ldots, \beta_m$, concretely
\begin{align*}
\sum_{k=0}^{m-1} c_k \cdot \phi^k \beta_1 = 1,
\quad 
\sum_{k=0}^{m-1} c_k \cdot \phi^k \beta_j = 0 \quad j=2,\ldots,n.
\end{align*}
This is exactly the first row vector of $M^{-1}$. 
%
In general, 
\[
\begin{pmatrix}
b_1(z)
\\
b_2(z)
\\
b_3(z)
\\
\vdots
\\
b_m(z)
\end{pmatrix}
= M^{-1}\cdot 
\begin{pmatrix}
v_0
\\
\phi v_1
\\
\phi^2 v_2
\\
\vdots
\\
\phi^{m-1} v_{m-1}
\end{pmatrix},
\]
where $v_k = b(\phi^{-k} z)$.
% and $M\cdot c$ is the Galois orbit of $c = (c_1,\ldots, c_m)$ in coordinates of the basis.
% The rows of $M^{-1}$ give us the linear functionals that recover the $i$-th coordinate of $z$ from its Galois-orbit.
\ulrich{%
\small
Invertibility of $M$ follows from that of $M^t$: 
any linear functional $(c_1,\ldots, c_m)$ in the kernel of $M^t$ gives rise to the linearizable polynomial
\[
p(X) = \sum_{i=0}^{m-1} c_i\cdot X^{2^i} 
\]
which has $\beta_1,\ldots, \beta_m$, and hence the whole space $F$ as roots. 
By dimension, $p(X)$ must be trivial, meaning that all its coefficents are zero.
}




